% Responsável: P2
\section{Configuração experimental}\label{sec:config}

\paragraph{Divisão desenvolvimento/teste.}
Antes de qualquer modelo, separámos 20\% dos dados para teste. A unidade de divisão é o
\emph{grupo}: o par de notícias, juntando os pares que partilham um texto ou uma ligação
(18 grupos com mais de um par). Usamos \texttt{StratifiedGroupKFold} (semente 42), estratificando
pelo tema. Resultam 5\,756 documentos de desenvolvimento e 1\,444 de teste, com zero grupos
partilhados. Sem este agrupamento, a notícia verdadeira de um par poderia ficar no treino e a
falsa no teste, e o modelo aproveitaria o assunto conhecido~\cite{kapoor2023}. A divisão está
guardada em \texttt{data/split.json} e é igual para todos os modelos.

\paragraph{Validação cruzada e métricas.}
Todas as escolhas (representação, hiperparâmetros, número de tópicos $K$ e passos de pré-processamento) foram tomadas com
validação cruzada de 5 folds, também agrupada, no desenvolvimento. Reportamos F1-macro (média e desvio-padrão entre folds), F1 por classe e o intervalo de confiança
a 95\% calculado com 1\,000 reamostragens \emph{bootstrap} das previsões \emph{out-of-fold}.
O teste é usado uma única vez, no fim,
com as configurações fixadas, e aí reportamos também a matriz de confusão.

\paragraph{Generalização por tema.}
Num protocolo \emph{leave-one-topic-out} dentro do desenvolvimento, treinamos sem um dos quatro
temas e avaliamos nesse tema.

\paragraph{Hiperparâmetros.}
Regressão logística com $C=10$; SVM linear com $C=1$; TF-IDF com $\mathit{min\_df}=2$ (palavras)
e 3 (caracteres), até 200\,000 e 300\,000 atributos; LDA com 30 iterações em lote.
Estes valores servem para comparar representações na Tabela~\ref{tab:cv}. Afinámos depois, por CV no
dev e na versão cortada, os três modelos clássicos mais fortes (39 configurações, Tabela~\ref{tab:tuning}):
em caso de empate (diferença $<0{,}001$) escolhemos o $C$ menor, isto é, o modelo mais regularizado.
Os ganhos são pequenos (até 0,6 pontos): caracteres com SVM passam de 0,934 para 0,939 com n-gramas de
1 a 4 e $C=0{,}3$, e palavras com SVM de 0,921 para 0,927 sem conversão para minúsculas, o que confirma
a \emph{ablation}. Os valores de CV das configurações escolhidas são ligeiramente otimistas, por terem
servido para escolher; o teste dá a estimativa sem viés.

\begin{table}[t]\centering\scriptsize
\caption{Afinação de hiperparâmetros por CV (5 folds, dev, texto cortado): número de configurações, F1-macro da pior e da melhor, e parâmetros escolhidos.}\label{tab:tuning}
\begin{tabular}{lrccl}
\toprule
Modelo & Config. & Pior & Melhor & Escolhida \\
\midrule
TF-IDF palavras + SVM & 24 & 0,891 & 0,927 & pipeline=raw, min\_df=2, C=1 \\
TF-IDF caracteres + SVM & 12 & 0,919 & 0,939 & n-gramas=(1, 4), min\_df=3, C=0.3 \\
TF-IDF caracteres + LogReg & 3 & 0,927 & 0,935 & n-gramas=(2, 5), min\_df=3, C=10 \\
\bottomrule
\end{tabular}
\end{table}


\paragraph{Ambiente e reprodutibilidade.}
Python 3.13 (testado também em 3.14 no Windows), versões fixadas em \texttt{requirements.txt},
semente 42 em todo o código. O comando \texttt{python run\_all.py} reproduz todas as tabelas;
o BERTimbau corre num \emph{notebook} no Google Colab (GPU T4) com a mesma divisão e os mesmos folds. Verificámos a reprodutibilidade clonando o repositório
público numa máquina limpa e correndo esse comando: todas as tabelas de resultados coincidem até
à última casa decimal, à exceção dos tempos de execução.
